% Generated draft: requires mathematical and source review.
\documentclass[11pt]{article}
\usepackage[margin=25mm]{geometry}
\usepackage{fontspec}
\setmainfont{Latin Modern Roman}
\usepackage{amsmath,amssymb,mathtools,unicode-math}
\setmathfont{Latin Modern Math}
\usepackage{enumitem,needspace,xurl,hyperref}
\setlength{\parindent}{0pt}
\setlength{\parskip}{5pt}
\setlength{\emergencystretch}{3em}
\newcommand{\N}{\mathbb{N}}
\newcommand{\Z}{\mathbb{Z}}
\newcommand{\Q}{\mathbb{Q}}
\newcommand{\R}{\mathbb{R}}
\newcommand{\C}{\mathbb{C}}
\newcommand{\F}{\mathbb{F}}
\newcommand{\A}{\mathbb{A}}
\newcommand{\PP}{\mathbb{P}}
\newcommand{\E}{\mathbb{E}}
\newcommand{\eps}{\varepsilon}
\newcommand{\dd}{\,\mathrm{d}}
\begin{document}
\section*{UnsolvedMath ID: 20002708}
\textbf{Source problem code:} catalogue entry
\subsection{Definitions and mathematical statement}
\textbf{Underlying source question.} Extend the type-\(A\), \(k=1\), Garsia--Haiman \((q,t)\)-Catalan/diagonal-harmonic connection to other reflection groups \(W\). The source compares earlier ``numbers (1)'' with the specialization \(t=q^{-1}\), up to a power of \(q\), but supplies neither those numbers nor the shift. Its verification identifies the actual heading as Problem 2.2, not the extracted 2.1. The source code and permanent identity are retained.

For a complexified reflection representation \(V\), define the diagonal coinvariant ring
\[
 D_W=\C[V\oplus V]/\langle\C[V\oplus V]^W_+\rangle,
\]
with the degrees in the two copies as its bigrading. This is the quotient displayed in the source, which calls it diagonal harmonics; dual harmonic-space conventions require care. The question seeks an appropriate quotient, representation and determinant/sign-isotypic bigrading. It records Haiman's observation that the naive finite-torus comparison does not persist as stated in type \(B\), and Gordon's further quotient as positive progress. The original finite-torus sentence suppresses a possible determinant twist, which must not be erased from the aid's explicit formula.

\textbf{Reported type-\(B_n\) Gordon quotient character.} Let \(W=B_n\) be the signed permutation group on \(V=\C^n\), let \(Q=\Z^n\) be its root lattice and \(h=2n\). For a signed cycle of \(w\), call it positive if the product of its signs is \(+1\), and let \(c_+(w)\) count those cycles. On the finite torus \(Q/(2n+1)Q\), the fixed-point count is
\[
 \#\operatorname{Fix}(w)=(2n+1)^{c_+(w)}.
\]
The odd modulus makes negative-cycle fixed coordinates zero; no primality assumption is needed. For the reported Gordon quotient \(R\) of \(D_{B_n}\), the aid gives
\[
 \chi_R(w)=\det(w)(2n+1)^{c_+(w)},\qquad
 \dim R=(2n+1)^n.
\]
Thus the relevant representation is the determinant twist of the finite-torus permutation representation, not an unqualified untwisted equality. The precise construction and general hypotheses of Gordon's quotient theorem are not supplied here.

\textbf{Reported complete rank-two case.} For \(B_2\), the aid attributes dimension \(25\) of the full diagonal coinvariant ring to Stump's dihedral Hilbert series. Its Gordon quotient also has dimension \(25\), so the surjective quotient is an isomorphism in this case. Let \(\det\) be the determinant character, let \(\chi_x,\chi_y\) be the two other nontrivial one-dimensional characters, and let \(\rho\) be the two-dimensional reflection representation. One can label \(\chi_x\) as the underlying-permutation sign and \(\chi_y\) as the product of coordinate signs; interchanging the two labels does not change the formula. The aid reports
\[
 D_{B_2}\cong\mathbf1\oplus6\det\oplus3\chi_x\oplus3\chi_y\oplus6\rho.
\]
Its determinant-isotypic multiplicity Hilbert series is
\[
 q^4+q^3t+q^2t^2+qt^3+t^4+qt.
\]
At \(q=t=1\) this has multiplicity six, consistently with the ungraded decomposition; the full dimensions sum to \(25\). No claim that the full ring equals Gordon's quotient in every type \(B_n\), or that this answers every \((q,t)\)-Catalan question for all \(W\), is made. In noncrystallographic settings a root-lattice finite torus is itself an additional construction issue. The omitted earlier formula and grading shift prevent an exact reconstruction of the original specialization target.
\subsection{Short English statement}
The aid gives an explicit determinant-twisted finite-torus character for Gordon's type-\(B\) quotient and a complete \(B_2\) decomposition/bigrading. It distinguishes that quotient from the full ring in general and retains the missing source formula, grading shift and problem-number correction.
\subsection{Sources}
\begin{itemize}
\item Local UnsolvedMath registry, permanent ID \(20002708\), catalogue entry, \texttt{aim-probability-notes.json}, \texttt{notes[149]}, workshop \emph{Braid groups, clusters and free probability}: \url{https://aimath.org/WWN/braidgroups/braidgroups.pdf}. Supplied verification identifies Problem 2.2; the extracted/source-item field says 2.1.
\item Source references: Garsia--Haiman \((q,t)\)-Catalan theory and Mark Haiman [29]; Iain Gordon's further quotient [26]. The bibliography entries are not supplied.
\item Supplied research summary cites Stump's dihedral Hilbert series and gives the signed-cycle quotient character and \(B_2\) data without full primary references. Attempt 1: partial\_result, partially\_solved, rigor ``proved'', low novelty confidence. The machine-generated aid is not peer reviewed and requires independent verification before citation; no PDF visit is claimed.
\end{itemize}
\end{document}
